% Generated by `taskgen course examples` from number_theory.linear_congruence (seed 6).
% Do not edit: change examples.yaml and regenerate.
\begin{example}[линейное сравнение с $d > 1$]\label{ex:euclid-linear-congruence}
Решить сравнение \[129 x \equiv 207 \pmod{381},\] т.\,е. найти все его решения $0 \le x < 381$ (или показать, что решений нет).

\par\noindent\textit{Решение.} Найдём $d = \gcd(a, n)$ и $v$ с $u n + v a = d$ расширенным алгоритмом Евклида для $(n, a) = (381, 129)$:\[\begin{array}{r|r|r|r}
r_i & q_i & u_i & v_i \\ \hline
381 &  & 1 & 0 \\
129 &  & 0 & 1 \\
123 & 2 & 1 & -2 \\
6 & 1 & -1 & 3 \\
3 & 20 & 21 & -62 \\
0 & 2 & -43 & 127 \\
\end{array}\qquad r_i = u_i \cdot 381 + v_i \cdot 129\]
$d = 3$, $v = -62$, т.\,е. $129 \cdot -62 \equiv 3\ (\textrm{mod}\ 381)$. Так как $d \mid b$ ($b/d = 69$), сравнение равносильно $43 x \equiv 69 \pmod{127}$, $n' = n/d = 127$. Умножая $a v \equiv d \pmod n$ на $b/d$: $x_0 \equiv v \cdot \frac{b}{d} = -62 \cdot 69 \equiv 40 \pmod{127}$. По модулю $n$ получаем $d = 3$ решений $x_0 + k n'$, $k = 0, \dots, 2$.\par Проверка: $129 \cdot 40 = 5160 \equiv 207\ (\textrm{mod}\ 381)$.\[x \in \{40, 167, 294\}.\]
\end{example}
